%% TeXora template: Midnight. Dark slides with amber highlights, easy on the eyes in a dim room.
%% Change the colours on the \definecolor lines below; everything else follows them.
% !TEX program = pdflatex
\documentclass[aspectratio=169]{beamer}
\usetheme{default}
\usepackage{amsmath}
\definecolor{night}{RGB}{17,24,39}
\definecolor{nightsoft}{RGB}{31,41,62}
\definecolor{amber}{RGB}{251,176,59}
\definecolor{mist}{RGB}{226,232,240}
\setbeamercolor{background canvas}{bg=night}
\setbeamercolor{normal text}{fg=mist,bg=night}
\setbeamercolor{structure}{fg=amber}
\setbeamercolor{title}{fg=white}
\setbeamercolor{subtitle}{fg=amber}
\setbeamercolor{frametitle}{fg=amber}
\setbeamercolor{alerted text}{fg=amber}
\setbeamercolor{block title}{bg=amber,fg=night}
\setbeamercolor{block body}{bg=nightsoft,fg=mist}
\setbeamercolor{block title example}{bg=teal!70!white,fg=night}
\setbeamercolor{block body example}{bg=nightsoft,fg=mist}
\setbeamercolor{block title alerted}{bg=red!65!white,fg=night}
\setbeamercolor{block body alerted}{bg=nightsoft,fg=mist}
\setbeamercolor{section in toc}{fg=mist}
\setbeamerfont{title}{series=\bfseries,size=\LARGE}
\setbeamerfont{frametitle}{series=\bfseries}
\setbeamertemplate{blocks}[rounded]
\setbeamertemplate{itemize items}[circle]
\setbeamertemplate{navigation symbols}{}
\setbeamertemplate{footline}{\hfill{\usebeamercolor[fg]{structure}\footnotesize\insertframenumber\,/\,\inserttotalframenumber}\hspace{1.4em}\vspace{0.9em}}
\setbeameroption{show notes on second screen=right}

\title{Your Presentation Title}
\subtitle{A short subtitle}
\author{Your Name}
\institute{Your Institution}
\date{\today}

\begin{document}

\begin{frame}
  \titlepage
  \note{Welcome everyone. Say who you are and what this talk covers.}
\end{frame}

\begin{frame}{Outline}
  \tableofcontents
  \note{Give the plan in one sentence.}
\end{frame}

\section{Introduction}
\begin{frame}{Key Ideas}
  \begin{itemize}
    \item<1-> Your first point
    \item<2-> Your second point, shown on the next step
    \item<3-> Your third point
  \end{itemize}
  \vspace{0.6em}
  \begin{block}<3->{Takeaway}
    One sentence the audience should remember.
  \end{block}
  \note{Each point appears on its own step. Pause after the takeaway.}
\end{frame}

\section{Details}
\begin{frame}{Two Columns}
  \begin{columns}[T]
    \begin{column}{0.47\textwidth}
      \begin{exampleblock}{Example}
        A short example, a result or a figure.
      \end{exampleblock}
    \end{column}
    \begin{column}{0.47\textwidth}
      \begin{alertblock}{Watch out}
        A caveat or a common mistake.
      \end{alertblock}
    \end{column}
  \end{columns}
  \note{Compare the two sides, left to right.}
\end{frame}

\begin{frame}{An Equation}
  \begin{equation*}
    f(x) = \sum_{k=0}^{n} a_k\, x^k
  \end{equation*}
  Explain what each symbol means, and why it matters here.
  \note{Read the equation aloud once, slowly.}
\end{frame}

\section{Conclusion}
\begin{frame}{Summary}
  \begin{enumerate}
    \item Your first conclusion
    \item Your second conclusion
  \end{enumerate}
  \vfill
  \begin{center}
    {\Large Thank you. Questions?}
  \end{center}
  \note{Thank the audience and invite questions.}
\end{frame}

\end{document}
